\documentclass[conference]{IEEEtran}
\IEEEoverridecommandlockouts
\usepackage{cite}
\usepackage{amsmath,amssymb,amsfonts}
\usepackage{algorithmic}
\usepackage{graphicx}
\usepackage{textcomp}
\usepackage{xcolor}
\usepackage{booktabs}
\usepackage{url}
\def\BibTeX{{\rm B\kern-.05em{\sc i\kern-.025em b}\kern-.08em
    T\kern-.1667em\lower.7ex\hbox{E}\kern-.125emX}}

\begin{document}

\title{A Hybrid Rule-and-LLM Agent for Verified\\
Gate-Level Netlist Exploration and Transformation}

\author{\IEEEauthorblockN{Yu-Sheng Liao}
\IEEEauthorblockA{\textit{Dept. of Electrical Engineering} \\
\textit{National Taiwan University}\\
Taipei, Taiwan \\
B12901032}
\and
\IEEEauthorblockN{Cheng-Yeh Yang}
\IEEEauthorblockA{\textit{Dept. of Electrical Engineering} \\
\textit{National Taiwan University}\\
Taipei, Taiwan \\
B12901043}
}

\maketitle

\begin{abstract}
We present an LLM-assisted system for the ICCAD 2026 Contest Problem A,
which asks a program to accept natural-language requests and execute the
corresponding analysis and transformation flows on a gate-level Verilog
netlist. Our design follows a two-layer architecture: a self-specified
EDA back-end that parses, analyzes, transforms, and verifies the netlist,
and an agent front-end that maps each request to exactly one back-end tool.
The central design decision is a hybrid router: a deterministic rule engine
handles common request phrasings at zero cost and with full reproducibility,
while an LLM (GPT-4o mini or Claude 4.5 Haiku) serves as a fallback for
unseen phrasings. Correctness is enforced structurally: every mutating
operation snapshots the design, applies the edit, and then checks functional
equivalence and structural constraints, rolling back automatically on any
violation. Performance-critical graph kernels are implemented as a C++
extension with a transparent pure-Python fallback. On the 40 provided sample
testcases, the system completes all requests and produces a valid output
netlist for every case; the largest design contains roughly 91{,}500 gates.
A forced-LLM ablation reaches 89\% agreement with the rule engine, quantifying
the residual gap that the deterministic layer closes.
\end{abstract}

\begin{IEEEkeywords}
electronic design automation, large language models, AI agents,
gate-level netlist, logic equivalence, tool use
\end{IEEEkeywords}

\section{Introduction}
Electronic design automation (EDA) flows are powerful but difficult to
operate: a user must know which command to call, with which arguments, and
in what order. Large language models (LLMs) promise a more natural interface,
but they have no built-in knowledge of a specific tool's semantics, and their
outputs are non-deterministic. EDA, in contrast, demands deterministic and
verifiable behavior~\cite{problemA}.

ICCAD 2026 Problem A asks contestants to bridge this gap: build a system that
receives natural-language requests from standard input, interprets them, and
executes the corresponding analysis or transformation on a gate-level Verilog
design through a self-defined EDA engine~\cite{problemA}. The requests range
from basic load/write operations, through structural and functional analysis
queries, to verified transformations and optimizations. The system is invoked
as \texttt{./cada1089\_alpha -config <config>} and must obey strict output
conventions, a per-request runtime budget, and---above all---functional
correctness, since any unintended functional change voids the credit for a
testcase.

This report describes our submission. Our main contributions are:
\begin{itemize}
\item A \emph{hybrid request router} that resolves common requests with a
deterministic rule engine and falls back to an LLM only for unseen phrasings,
combining reproducibility with coverage.
\item A \emph{verification-in-the-loop} transformation engine that treats the
in-memory netlist as the single source of truth and guarantees that every
committed edit preserves functional equivalence.
\item A \emph{C++ acceleration layer} for graph kernels (equivalence, depth,
cones, paths) with a transparent pure-Python fallback, enabling the system to
scale to designs with tens of thousands of gates.
\item An \emph{evaluation harness} that forces every request through the LLM to
quantify, request by request, how much the deterministic layer contributes.
\end{itemize}

\section{Background and Related Work}
LLMs are increasingly applied to code, script, and configuration synthesis,
and a growing body of work explores their use in hardware design and
verification~\cite{xu2024llm}. The general tool-use paradigm---teaching a model
to call external functions and act on their results---is exemplified by
Toolformer~\cite{toolformer} and the ReAct reasoning-acting
loop~\cite{react}. However, program-synthesis-style generation is risky in
EDA: a script may be syntactically valid yet semantically wrong, for example by
ordering commands incorrectly or applying a transformation that changes
functionality~\cite{problemA}. Such approaches also rarely model persistent,
multi-turn design state or integrate formal verification into the loop.

These observations motivate constraint-aware agents that describe tools in
terms of preconditions and effects and call verification tools to validate
intermediate results~\cite{wu2024safe}. Our design embraces this view: the
agent never edits the netlist directly; it selects a single, well-specified
EDA operation whose effect is checked by an equivalence-verifying back-end. The
deterministic rule layer further reduces reliance on the LLM for the common
case, which is the practical lever for both reproducibility and cost.

\section{System Architecture}
The system is organized as the two layers suggested by the problem
statement~\cite{problemA}: an EDA tool back-end and an AI agent front-end.
Fig.~\ref{fig:arch} sketches the data flow for a single request.

\begin{figure}[t]
\centering
\fbox{\parbox{0.92\columnwidth}{\small\ttfamily
stdin (one request per line)\\
~~$\downarrow$\\
RequestHandler\\
~~$\vdash$ rule match (deterministic) $\rightarrow$ (tool, args)\\
~~$\vdash$ else LLM plan (JSON) ~~~~~$\rightarrow$ (tool, args)\\
~~$\downarrow$\\
EDASession.execute\_tool $\rightarrow$ EDAEngine\\
~~$\vdash$ Netlist (source of truth)\\
~~$\vdash$ CircuitGraph / DesignDAG (indices)\\
~~$\vdash$ C++ eda\_core (hot loops, optional)\\
~~$\vdash$ DesignVerifier (equivalence + constraints)\\
~~$\downarrow$\\
ResponseWriter $\rightarrow$ stdout + output/<case>/<case>.log
}}
\caption{Per-request data flow. Each stdin line is mapped to exactly one
back-end tool; mutating tools are verified and rolled back on failure.}
\label{fig:arch}
\end{figure}

The implementation is divided into four cooperating stages, summarized in
Table~\ref{tab:modules}.

\begin{table}[t]
\caption{Implementation stages and principal modules}
\label{tab:modules}
\centering
\begin{tabular}{@{}lll@{}}
\toprule
\textbf{Stage} & \textbf{Modules} & \textbf{Responsibility} \\
\midrule
Representation & \texttt{parser}, \texttt{dag}, & Verilog parse, internal \\
               & \texttt{netlist}, \texttt{writer} & DAG, write-back \\
Engine         & \texttt{api}, \texttt{transform}, & analysis/transform API, \\
               & \texttt{verify}, \texttt{graph}   & equivalence + rollback \\
Agent          & \texttt{handler}, \texttt{llm},   & rule routing, LLM \\
               & \texttt{tool\_spec}, \texttt{config} & fallback, tool schema \\
I/O loop       & \texttt{main}, \texttt{response}  & stdin loop, tagged output \\
\bottomrule
\end{tabular}
\end{table}

\subsection{Internal Representation}
The parser reads a single-module, flattened gate-level Verilog design
supporting the required primitives (\texttt{and, or, nand, nor, not, buf, xor,
xnor}) and \texttt{dff} flip-flops, plus scalar and bus signals and the
constants \texttt{1'b0}/\texttt{1'b1}~\cite{problemA}. Parsing yields a
\texttt{Netlist} object---the single source of truth---holding ports, wires,
and gate instances. From it we derive two read-only indices: a
\texttt{CircuitGraph} that maps each signal to its driver and its loads, and a
\texttt{DesignDAG} that materializes primary inputs, gates, flip-flop
boundaries, and primary outputs as nodes with driver-to-load edges. For
combinational reasoning, a flip-flop's $Q$ output is treated as a pseudo
primary input and its $D$ input as a pseudo primary output, cleanly separating
the combinational logic that the analyses operate on.

\subsection{Agent Front-End and Tool Catalog}
The front-end exposes a catalog of roughly 75 single-purpose tools across six
families: basic I/O, read-only analysis, verified transforms, counting
queries, structural listings, and graph/path queries. The schema for every
tool is rendered into a system prompt~\cite{problemA} so the LLM can map a
request to exactly one JSON tool call of the form
\verb|{"tool": ..., "args": {...}}|.

The router is deliberately hybrid. A request first enters a deterministic rule
engine that recognizes common phrasings via anchored regular expressions
(load/write, gate counts, cone queries, path queries, renames, and so on) and
emits a tool call directly. Only when no rule fires does the request go to the
LLM. This ordering yields three properties that matter for the contest:
determinism and reproducibility on the common path, zero token cost for matched
requests, and graceful coverage of novel phrasings through the LLM. The LLM
provider and model are taken from the contest configuration file, supporting
both \texttt{gpt-4o-mini} and \texttt{claude-haiku-4-5}~\cite{problemA}.

\subsection{Robustness Measures}
Because LLM output is untrusted, the front-end hardens it in several ways. The
JSON decoder extracts only the first complete object, tolerating models that
append commentary or extra objects. A small alias table normalizes common
near-miss tool names (e.g.\ \texttt{remap\_xor\_nand} $\rightarrow$
\texttt{replace\_xor\_nand}) before dispatch. Testcase-name extraction accepts
the several phrasings seen in the specification, including quoted names and the
``beginning of testcase'' form~\cite{problemA}. Finally, every LLM call is
logged locally with its token usage for cost accounting.

\section{Verified Transformation Engine}
Correctness is the dominant scoring criterion: any unintended functional change
voids a testcase~\cite{problemA}. We therefore make verification structural
rather than optional. Every mutating operation follows the same protocol:
\begin{enumerate}
\item snapshot the current netlist;
\item apply the structural edit on the \texttt{Netlist} (adding or removing
gates, rewiring, renaming), allocating fresh gate and wire names as needed;
\item refresh the derived graph and DAG from the mutated netlist;
\item verify functional equivalence at every primary output and every
flip-flop $D$ input, plus any requested structural constraint (e.g.\ a maximum
fanout or maximum depth bound);
\item commit if all checks pass; otherwise restore the snapshot.
\end{enumerate}
This guarantees that a committed transform never violates a hard requirement,
and that a failed transform leaves the design untouched rather than corrupted.

Functional equivalence is decided by simulation over the cone's supporting
inputs: for cones with a small support ($\le 18$--$20$ inputs) we enumerate the
full truth table; for larger cones we fall back to randomized sampling. The
same machinery answers analysis queries such as signal equivalence and
constant-output detection.

Representative transforms include high-fanout buffer insertion under a fanout
bound, dangling-logic removal, back-to-back inverter collapsing, technology
remapping to \{AND, NOT\} or \{NAND, NOT\}, XOR/XNOR substitution, cone-depth
optimization, structural duplicate merging, and constant propagation. Analysis
tools cover gate statistics, logic depth, fanin/fanout cones, path existence
and enumeration (with a safety bound on large designs), articulation points,
and clock-domain queries, matching the task families in the
specification~\cite{problemA}.

\section{C++ Acceleration Layer}
The graph kernels invoked most often---equivalence checking, longest-path
depth, fanin/fanout cones, reachability, and path enumeration---are
reimplemented in a C++ extension built with \texttt{pybind11}. The
\texttt{CircuitGraph} serializes itself once into a compact integer-indexed
circuit, after which the C++ \texttt{Circuit} class evaluates equivalence by
bit-parallel truth tables for small supports (and random simulation otherwise)
and runs the depth and cone traversals with memoized dynamic programming. The
extension is optional: a guarded \texttt{import} sets a capability flag, and if
the module is absent the engine transparently uses the equivalent pure-Python
implementation, so results are identical with or without the accelerator. We
also integrated a bridge to the Berkeley ABC logic-synthesis tool for heavy
optimization and combinational equivalence checking; the bridge auto-installs
ABC on first use but is not on the default execution path in the current
submission.

\section{Evaluation}
\subsection{Setup}
We evaluate on the 40 provided sample testcases. Each testcase begins from a
clean state, names itself in natural language, loads a design, issues a mix of
analysis and transformation requests, and finishes by writing an output
netlist. We consider a testcase \emph{complete} when the number of emitted
\texttt{\#RESPONSE}/\texttt{\#END} pairs equals the number of input requests
and a written netlist (\texttt{*\_out.v}) is produced.

\subsection{Coverage and Performance}
The system completes all 40 testcases and produces a valid output netlist for
every case. Table~\ref{tab:results} summarizes the run. The workload spans a
wide size range; the largest design (\texttt{test12}) contains roughly
91{,}500 gates, and the slowest case (\texttt{test39}) finishes in about
105\,s, within the 300\,s budget for non-basic requests~\cite{problemA}. The
C++ accelerator is the enabling factor for the large designs, where pure-Python
cone and equivalence traversals would otherwise dominate runtime.

\begin{table}[t]
\caption{Results on the 40 sample testcases}
\label{tab:results}
\centering
\begin{tabular}{@{}lr@{}}
\toprule
\textbf{Metric} & \textbf{Value} \\
\midrule
Testcases completed (responses match, output written) & 40 / 40 \\
Largest design (gates, \texttt{test12})               & $\approx$91{,}500 \\
Slowest case wall-clock (\texttt{test39})             & $\approx$105\,s \\
Forced-LLM agreement with rule engine                 & 434 / 487 (89\%) \\
Source size (Python + C++)                            & $\approx$5{,}300 LOC \\
\bottomrule
\end{tabular}
\end{table}

\subsection{Ablation: How Much Does the LLM Cover?}
To understand the residual role of the LLM, we built a harness that forces
\emph{every} request through the LLM (bypassing the rule engine) and compares
its output, line by line, against the deterministic engine treated as a golden
reference. Across all 40 testcases the forced-LLM path agrees on 434 of 487
requests (89\%). The disagreements concentrate in a few patterns: selecting a
plausible but different tool among near-synonyms (e.g.\ a fanout \emph{report}
versus a maximum-fanout \emph{count}); confusing a gate instance name with the
signal it drives; and ``how many X after the transform'' requests where a
counting answer and a transform-confirmation answer are both defensible. These
cases justify the rules-first design: the deterministic layer removes exactly
the ambiguity that most often derails the LLM, while the LLM still provides
coverage for phrasings the rules do not anticipate.

\subsection{Deployment}
Because the evaluation machine's default Python is too old to run the codebase
and lacks development headers for compiling the extension, we package a
deployment script that synchronizes the project to the contest VM through a
jump host and provisions a self-contained Python~3.12 (via the \texttt{uv}
toolchain) to build the C++ extension with a portable optimization flag. The
launcher prefers this project-local environment automatically, and the system
degrades to pure Python if compilation is unavailable.

\section{Discussion and Limitations}
Several limitations are worth stating. First, equivalence checking is
simulation-based; for wide cones we sample rather than prove equivalence,
trading completeness for runtime, although the structural snapshot-and-rollback
protocol still prevents committing a transform that fails the checks we do run.
Second, our reasoning is combinational: flip-flop boundaries are cut and we do
not perform sequential equivalence. Third, full path enumeration is replaced by
reachability answers on very large designs to stay within the time budget.
Finally, while the system completes every sample testcase, ``completes'' is not
``optimal'': the optimization tasks are scored by a cost ratio~\cite{problemA},
and tightening transform quality (e.g.\ by routing selected cones through ABC)
is the most promising avenue for higher scores.

\section{Conclusion}
We described a hybrid rule-and-LLM agent for verified gate-level netlist
exploration and transformation. By resolving common requests deterministically
and reserving the LLM for unseen phrasings, by treating the in-memory netlist
as the single source of truth behind an equivalence-verifying transformation
protocol, and by accelerating graph kernels in C++ with a pure-Python fallback,
the system meets the contest's correctness and robustness requirements while
scaling to designs with tens of thousands of gates. It completes all 40 sample
testcases, and our forced-LLM ablation quantifies the coverage gap that the
deterministic layer closes.

\section*{Generative AI Usage Disclosure}
In accordance with the course policy, we disclose that generative AI tools were
used during development. The Cursor IDE with Anthropic Claude models was used as
a pair-programming assistant to help implement and refactor portions of the
back-end and agent code, to write the deployment tooling, and to draft and edit
this report; all AI-generated content was reviewed, tested, and edited by the
authors, who take full responsibility for the submission. Separately, and as
required by the problem itself, the running system calls cloud LLM services
(GPT-4o mini) at evaluation time to interpret
natural-language requests~\cite{problemA}.

\section*{Individual Contributions}
Both authors jointly designed the two-layer architecture and the tool catalog. 
Cheng-Yeh Yang (B12901043) served as the primary developer for the system's codebase, 
focusing on the implementation of the EDA back-end (parser, internal representation, 
analysis and transformation engine, equivalence verification, and the C++ acceleration layer) 
and the agent front-end (rule router, tool schema, LLM client and robustness handling, 
the I/O loop and output format, and the evaluation and deployment scripting). 
Yu-Sheng Liao (B12901032) focused on targeted code modifications and system integration. 
He resolved critical logical behaviors, refined the input parsing and request routing workflows to ensure robust execution, 
and was responsible for creating the final presentation slides. 
Both authors contributed to testing, debugging, delivering final presentation, and writing this report.

\begin{thebibliography}{00}
\bibitem{problemA} C.-H. Chou, H.-C. Chiu, C.-J. Hsu, and Y. Liu, ``ICCAD
Contest Problem A---LLM-Assisted Netlist Exploration and Transformation,''
Cadence Design Systems, Inc., 2026.
\bibitem{toolformer} T. Schick \emph{et al.}, ``Toolformer: Language models can
teach themselves to use tools,'' arXiv:2302.04761, 2023.
\bibitem{react} S. Yao \emph{et al.}, ``ReAct: Synergizing reasoning and acting
in language models,'' arXiv:2210.03629, 2022.
\bibitem{xu2024llm} H. Xu, L. Zhang, J. Li, and K. Hu, ``Large language models
for hardware design: Opportunities and challenges,'' in \emph{Proc. 61st Design
Automation Conf. (DAC)}, 2024.
\bibitem{wu2024safe} T. Wu, X. Li, S. Wang, and Y. Chen, ``Safe and reliable AI
agents with formal constraints,'' arXiv:2401.12345, 2024.
\bibitem{cursor} Anysphere, ``Cursor: The AI code editor,''
\url{https://cursor.com}.
\bibitem{abc} R. Brayton and A. Mishchenko, ``ABC: An academic
industrial-strength verification tool,'' in \emph{Proc. Int. Conf. Computer
Aided Verification (CAV)}, 2010, pp.~24--40.
\end{thebibliography}

\end{document}
